
% =====================================================================
\section{The Vorton-Scalar}
\label{sec:vs}
% =====================================================================

This section meets objective O3. It invents the boson that mediates the pairing, defines it by a
Lagrangian, and derives its neutral current, its coupling to the fermions and its field equations.

\subsection{What is wanted}

Pairing needs an attraction between two fermions. Hypothesis H1 prescribes its origin: a \emph{scalar}
boson field, and an interaction of the kind known from the weak force, in which a \emph{neutral
current}, a current that does not change the kind of fermion on which it acts, couples to a current of
the fermions. The neutral current is to be built from the scalar alone.

A real scalar field has no current. The simplest field that has one is a \emph{complex} scalar
$\PhiVS$, whose phase rotations $\PhiVS\to e^{i\epsilon}\PhiVS$ are a symmetry; the conserved current
of that symmetry is bilinear in $\PhiVS$, and it is the neutral current sought. For it to act on the
fermions, the fermions must carry the same symmetry. A fermion field with a mass term has only one
phase symmetry, the fermion number. The massless field has a second one, the chiral rotation, and the
mass term breaks it. This suggests the construction: \emph{let the mass of the fermion field come
from the scalar}, so that the chiral rotation of the fermions, accompanied by a phase rotation of the
scalar, is an exact symmetry.

%%<1671>
\subsection{Definition}
\label{sec:vsdef}

The Vorton-Scalar is a complex scalar field $\PhiVS(x)$ on the $7+1$ spacetime. Together with the
fermion field it is governed by the Lagrangian
\begin{align}
\mathcal L_{1671}^{\rm VS}={}&\bar\psi\,i\gamma^A\partial_A\psi
+\eta^{AB}\,\partial_A\PhiVS^{*}\,\partial_B\PhiVS\nonumber\\
&-\lamVS\Big(|\PhiVS|^2-\frac{\vVS^2}{2}\Big)^{2}\nonumber\\
&-\yVS\,\bar\psi\big(\PhiVS\,P^{R}+\PhiVS^{*}\,P^{L}\big)\psi .
\label{eq:LVS}
\end{align}
Here $\psi=\psi_{1671}$, $\gamma^A=\gamma_{1671}^A$, $\eta=\eta_{1671}$, $\bar\psi=\psi^\dagger\beta$
with $\beta=\gamma^4$, and $P^{L,R}=\tfrac12(1\mp\chi)$ are the complex chiral projectors
(\ref{eq:PLPR}) with the chirality matrix $\chi=\chi_{1671}$. The three constants are real and
positive: the self-coupling $\lamVS$, the vacuum value $\vVS$, and the Yukawa coupling $\yVS$. The last
term is Hermitian, because $(\bar\psi P^{R}\psi)^\dagger=\psi^\dagger P^{R}\beta\psi=\bar\psi P^{L}\psi$:
the matrix $\chi$ anticommutes with $\beta$.

\emph{Symmetries.} Besides the Lorentz group of the $7+1$ spacetime, the Lagrangian (\ref{eq:LVS}) has
two phase symmetries,
\begin{align}
U(1)_F:&\quad\psi\to e^{i\epsilon}\psi;\nonumber\\
\UVS:&\quad\psi\to e^{i\epsilon\chi}\psi,\quad\PhiVS\to e^{-2i\epsilon}\PhiVS .
\label{eq:U1}
\end{align}
The first is the fermion number. The second is the symmetry wanted. To verify it, note that $\chi$ is
Hermitian and anticommutes with every $\gamma^A$, so that $\bar\psi\to\bar\psi e^{i\epsilon\chi}$ and
$\bar\psi\gamma^A\partial_A\psi\to\bar\psi e^{i\epsilon\chi}\gamma^Ae^{i\epsilon\chi}\partial_A\psi
=\bar\psi\gamma^A\partial_A\psi$; and
$\bar\psi P^{R}\psi\to\bar\psi e^{2i\epsilon\chi}P^{R}\psi=e^{2i\epsilon}\bar\psi P^{R}\psi$, which
the phase of $\PhiVS$ compensates. A bare mass term $m\bar\psi\psi$ is not invariant, and it is
therefore absent from Eq.~(\ref{eq:LVS}).
%%</1671>
%%<1644>
\subsection{The covariant theory has no chiral phase symmetry}
\label{sec:nogo}

The construction just described fails for the Lagrangian (\ref{eq:lagrangian}) that is invariant
under the whole group $\Spin_0(4,4)$. On the $4+4$ spacetime the chirality matrix $\omega$ is real,
and $\Cm$, a product of four Dirac matrices, \emph{commutes} with it. Under the chiral phase rotation
$\psi\to e^{i\epsilon\omega}\psi$ one finds, with $\tilde\psi=\psi^\dagger\Cm$,
\begin{equation}
\tilde\psi\psi\to\tilde\psi\psi,\qquad
\tilde\psi\gamma^A\partial_A\psi\to\tilde\psi\,e^{-2i\epsilon\omega}\gamma^A\partial_A\psi .
\end{equation}
The \emph{mass term} is invariant and the \emph{kinetic term} is not, the opposite of the situation in
ordinary four-dimensional theory and on the $7+1$ spacetime. The kinetic term is invariant instead
under the real, non-compact rescaling $\psi\to e^{\epsilon\omega}\psi$, which multiplies the two
8-component chiral parts by $e^{-\epsilon}$ and $e^{+\epsilon}$ and mixes $\tilde\psi\psi$ with
$\tilde\psi\omega\psi$ by a hyperbolic rotation. A scalar that compensated this rotation would have
the invariant $(\mathrm{Re}\,\PhiVS)^2-(\mathrm{Im}\,\PhiVS)^2$, of indefinite sign: one of its two
components would have negative kinetic energy.

\emph{In the fully covariant $4+4$ theory the construction does not exist.} This is a consequence of
the real chirality of the $4+4$ signature, and it is recorded here as a result.

\subsection{Definition in the quantized sector}
\label{sec:vsdef}

Section~\ref{sec:quantized} defined the quantized field by the restriction to fields that do not
depend on the additional times, with the Lagrangian (\ref{eq:Lred}). That Lagrangian differs from the
covariant one by the grading $W$, and $W$, which contains an odd number of Dirac matrices,
\emph{anticommutes} with $\omega$. In $\mathcal L^{\rm red}$ the roles are therefore the familiar
ones: the term $\psi^\dagger(i\partial_4+i\alpha^j\partial_j)\psi$ is invariant under
$\psi\to e^{i\epsilon\omega}\psi$, because $\omega$ commutes with the $\alpha^j$, and the mass term
$m\,\psi^\dagger\beta\psi$ is not, because $\omega$ anticommutes with $\beta$.

In the formulas that this paper shares with its companion the chirality matrix is denoted by $\chi$
and the adjoint spinor of the quantized theory by $\bar\psi$. For the present spacetime
\begin{equation}
\chi\equiv\omega_{1644},\qquad\bar\psi\equiv\psi^\dagger\beta_{1644},\qquad P^{L,R}=\tfrac12(1\mp\chi).
\label{eq:chi}
\end{equation}
The matrix $\chi$ is real, and $\bar\psi$ must not be confused with the covariant adjoint
$\tilde\psi=\psi^\dagger\Cm$ of Sec.~\ref{sec:dirac}.

The Vorton-Scalar of the $4+4$ model is a complex scalar field $\PhiVS$ that, like the fermion field,
does not depend on $x5,x6,x7$. It is governed by the Lagrangian
\begin{align}
\mathcal L_{1644}^{\rm VS}={}&\psi^\dagger\big(i\partial_4+i\alpha^j\partial_j\big)\psi
+|\partial_4\PhiVS|^2-\sum_{j=0}^{3}|\partial_j\PhiVS|^2\nonumber\\
&-\lamVS\Big(|\PhiVS|^2-\frac{\vVS^2}{2}\Big)^{2}\nonumber\\
&-\yVS\,\bar\psi\big(\PhiVS\,P^{R}+\PhiVS^{*}\,P^{L}\big)\psi ,
\label{eq:LVS}
\end{align}
with $\psi=\psi_{1644}$, $\alpha^j=\alpha_{1644}^j$ and $j$ running over the four space coordinates.
The three constants are real and positive: the self-coupling $\lamVS$, the vacuum value $\vVS$, and
the Yukawa coupling $\yVS$. The Lagrangian is invariant under $\Spin(4,1)\times\Spin(3)$, the
rotations and boosts of $(x0,x1,x2,x3;x4)$ and the rotations of the three additional times, and under
two phase symmetries,
\begin{align}
U(1)_F:&\quad\psi\to e^{i\epsilon}\psi;\nonumber\\
\UVS:&\quad\psi\to e^{i\epsilon\chi}\psi,\quad\PhiVS\to e^{-2i\epsilon}\PhiVS .
\label{eq:U1}
\end{align}
The first is the fermion number. The second is the symmetry wanted:
$\bar\psi P^{R}\psi\to e^{2i\epsilon}\bar\psi P^{R}\psi$, which the phase of $\PhiVS$ compensates. A
bare mass term is not invariant and is therefore absent.
%%</1644>

\emph{The name.} In the earlier study \cite{NashNCSS} the name Vorton-Scalar denoted the scalar that
mediates the pairing and that every quantized vortex of the superfluid carries along with it. That
meaning is kept (Sec.~\ref{sec:vortons}). The field invented here is in addition \emph{complex}, with
a broken phase symmetry. Such a field has line-like defects of its own, on which the fermions are
massless and can carry a persistent current, and closed loops stabilized by such a current are what
the literature on cosmic strings calls \emph{vortons} \cite{DavisShellard1989}.

\subsection{Dimensions of the couplings, and the reduction to four dimensions}
\label{sec:reduction}

In natural units the action is dimensionless. In eight dimensions the kinetic terms then give the
fields the dimensions $[\psi]=\text{mass}^{7/2}$ and $[\PhiVS]=\text{mass}^{3}$, hence
\begin{equation}
[\vVS]=\text{mass}^{3},\qquad[\yVS]=\text{mass}^{-2},\qquad[\lamVS]=\text{mass}^{-4}.
\end{equation}
Couplings of negative mass dimension mean that the model is an \emph{effective} theory, valid below
some scale. This is unavoidable for interacting fields in eight dimensions and is not specific to the
present model.

In the first stage of the program the fields that matter do not depend on the four extra coordinates:
the hidden coordinate is curled up, and the three others are assumed to play no dynamical role
\SIG{(for the additional times this is the restriction of Sec.~\ref{sec:quantized})}{}. If the extra
coordinates span a volume $V_4$, a field that is constant along them is
$\psi=\psi^{(4)}/\sqrt{V_4}$, $\PhiVS=\PhiVS^{(4)}/\sqrt{V_4}$, and the action becomes a
four-dimensional one with the dimensionless couplings and the scale
\begin{equation}
\yVS^{(4)}=\frac{\yVS}{\sqrt{V_4}},\qquad\lamVS^{(4)}=\frac{\lamVS}{V_4},\qquad\vVS^{(4)}=\vVS\sqrt{V_4},
\end{equation}
where $\vVS^{(4)}$ is a mass. From here on all couplings are these four-dimensional ones, and the
superscript is dropped. We call the fields that do not depend on the four extra coordinates the
\emph{zero-mode sector}.

\subsection{Symmetry breaking: the mass of the fermions and the two modes}
\label{sec:ssb}

The potential in Eq.~(\ref{eq:LVS}) is minimal at $|\PhiVS|=\vVS/\sqrt2$. Write
\begin{equation}
\PhiVS(x)=\frac{\vVS+\sVS(x)}{\sqrt2}\,e^{i\thVS(x)/\vVS}
\label{eq:polar}
\end{equation}
with two real fields, the \emph{radial mode} $\sVS$ and the \emph{phase mode} $\thVS$. Because
$\PhiVS P^{R}+\PhiVS^*P^{L}=|\PhiVS|\,e^{i\chi\arg\PhiVS}$, the Yukawa term becomes
$-m\,(1+\sVS/\vVS)\,\bar\psi\,e^{i\chi\thVS/\vVS}\psi$, and to first order in the boson fields
\begin{align}
\mathcal L_{\tg}^{\rm VS}={}&\mathcal L_{\psi}(m)+\tfrac12(\partial\sVS)^2-\tfrac12m_s^2\sVS^2+\tfrac12(\partial\thVS)^2\nonumber\\
&-\gVS\,\sVS\,\bar\psi\psi-\gVS\,\thVS\,\bar\psi\,i\chi\,\psi+\dots,
\label{eq:Lbroken}
\end{align}
where $\mathcal L_{\psi}(m)$ is the free Lagrangian of the fermion field with the mass $m$,
\SIG{Eq.~(\ref{eq:Lred})}{Eq.~(\ref{eq:lagrangian})}, the dots stand for terms with three or more boson
fields, and
\begin{align}
m_{\tg}&=\frac{\yVS\,\vVS}{\sqrt2},\qquad m_s=\sqrt{2\lamVS}\;\vVS,\nonumber\\
\gVS&=\frac{m_{\tg}}{\vVS}=\frac{\yVS}{\sqrt2}.
\label{eq:masses}
\end{align}
Three facts follow. First, the fermion mass is no longer a free parameter put in by hand: it is the
Yukawa coupling times the vacuum value of the Vorton-Scalar. Second, the radial mode is a scalar of
mass $m_s$ that couples to the scalar density $\bar\psi\psi$. Third, the phase mode is massless, being
the Goldstone boson \cite{Goldstone1961} of the broken symmetry $\UVS$, and it couples to the
pseudoscalar density $\bar\psi\,i\chi\,\psi$. Both couplings have the same strength
$\gVS=m/\vVS$: heavier fermions couple more strongly, a relation of the Goldberger-Treiman type
\cite{GoldbergerTreiman1958}. We also use
\begin{equation}
\aVS\equiv\frac{\gVS^2}{4\pi}.
\label{eq:alphaVS}
\end{equation}

In the Hamiltonian form of the theory, Eq.~(\ref{eq:schroedinger}), the two vertices read
\begin{align}
\text{radial mode:}&\quad-\gVS\,\sVS\,\psi^\dagger\beta\,\psi,\nonumber\\
\text{phase mode:}&\quad-\gVS\,\thVS\,\psi^\dagger(i\beta\chi)\psi .
\label{eq:vertices}
\end{align}
The matrices $\beta$ and $i\beta\chi$ are Hermitian. Equations (\ref{eq:schroedinger}),
(\ref{eq:hamiltonian}) and (\ref{eq:vertices}) have the same form for both spacetimes; all that
differs is which matrices $\beta$, $\alpha^j$ and $\chi$ stand in them. Everything from
Sec.~\ref{sec:pairing} on is derived from these three equations.

\subsection{The neutral current}
\label{sec:nc}

The conserved current of the symmetry $\UVS$ follows from Eq.~(\ref{eq:U1}) by Noether's theorem. It
is the sum of a current of the fermions and a current of the scalar. With the components along the
time $x4$ and along the space coordinates $xj$ of the model
\SIG{($j=0,1,2,3$)}{(the seven values $j\neq4$)},
\begin{align}
\JVS^4&=i\big(\PhiVS^{*}\,\partial_4\PhiVS-\PhiVS\,\partial_4\PhiVS^{*}\big),\nonumber\\
\JVS^j&=-i\big(\PhiVS^{*}\,\partial_j\PhiVS-\PhiVS\,\partial_j\PhiVS^{*}\big),\nonumber\\
J_5^4&=\psi^\dagger\chi\,\psi,\qquad J_5^j=\psi^\dagger\alpha^j\chi\,\psi,
\label{eq:currents}
\end{align}
the conservation law is
\begin{equation}
\partial_4\big(J_5^4-2\JVS^4\big)+\partial_j\big(J_5^j-2\JVS^j\big)=0 .
\label{eq:conservation}
\end{equation}
$\JVS$ is \emph{the neutral current of the Vorton-Scalar}: it is built from $\PhiVS$ alone, and it is
neutral in the sense of the weak interaction, since it leaves the fermion on which it acts unchanged.
$J_5$ is the axial current of the fermions. Neither is conserved separately: the current of the
Vorton-Scalar supplies the divergence of the axial current.
%%<1671>
In covariant form $\JVS^A=i(\PhiVS^{*}\partial^A\PhiVS-\PhiVS\partial^A\PhiVS^{*})$ and
$J_5^A=\bar\psi\gamma^A\chi\psi$, with the index raised by $\eta_{1671}$.
%%</1671>

The coupling of the phase mode in Eq.~(\ref{eq:Lbroken}) can be written with these currents. Redefine
the fermion field by a local chiral rotation, $\psi\to e^{-i\chi\thVS/2\vVS}\psi$. This removes
$\thVS$ from the Yukawa term, and the kinetic term of the fermions produces instead
$(1/2\vVS)(\partial_4\thVS\,J_5^4+\partial_j\thVS\,J_5^j)$. With Eq.~(\ref{eq:polar}),
$\JVS^4=-(\vVS+\sVS)^2\partial_4\thVS/\vVS$ and $\JVS^j=+(\vVS+\sVS)^2\partial_j\thVS/\vVS$, so that to
lowest order in the fields
\begin{equation}
\mathcal L_{\rm NC}=-\frac{\GVS}{\sqrt2}\,\big(\JVS^4J_5^4-\JVS^jJ_5^j\big),\qquad\frac{\GVS}{\sqrt2}=\frac{1}{2\,\vVS^{2}}.
\label{eq:LNC}
\end{equation}
This is a \emph{current-current interaction of the weak type}, with a constant $\GVS$ of the Fermi type
that is fixed by the vacuum value alone, just as the Fermi constant of the weak interaction is
$G_F/\sqrt2=1/(2v_{\rm EW}^2)$ in terms of the electroweak vacuum value $v_{\rm EW}$
\cite{Weinberg1967}. Equations (\ref{eq:Lbroken}) and (\ref{eq:LNC}) describe the same physics in two
sets of field variables. Equation (\ref{eq:Lbroken}) is used for the calculations below;
Eq.~(\ref{eq:LNC}) shows that the exchange of the phase mode \emph{is} the neutral-current interaction
that hypothesis H1 asks for.

\subsection{Field equations}

Varying Eq.~(\ref{eq:LVS}) with respect to $\bar\psi$ and $\PhiVS^*$ gives the coupled field equations
of the model,
%%<1671>
\begin{align}
&\big[i\gamma_{1671}^A\partial_A-\yVS\big(\PhiVS P_{1671}^{R}+\PhiVS^{*}P_{1671}^{L}\big)\big]\psi_{1671}=0,\nonumber\\
&\eta_{1671}^{AB}\partial_A\partial_B\PhiVS+2\lamVS\Big(|\PhiVS|^2-\frac{\vVS^2}{2}\Big)\PhiVS\nonumber\\
&\qquad\qquad+\yVS\,\bar\psi_{1671}P_{1671}^{L}\psi_{1671}=0 .
\label{eq:fieldeqs}
\end{align}
%%</1671>
%%<1644>
\begin{align}
&\big(i\partial_4+i\alpha_{1644}^j\partial_j\big)\psi_{1644}\nonumber\\
&\qquad-\yVS\,\beta_{1644}\big(\PhiVS P_{1644}^{R}+\PhiVS^{*}P_{1644}^{L}\big)\psi_{1644}=0,\nonumber\\
&\Big(\partial_4^2-\sum_{j=0}^{3}\partial_j^2\Big)\PhiVS+2\lamVS\Big(|\PhiVS|^2-\frac{\vVS^2}{2}\Big)\PhiVS\nonumber\\
&\qquad\qquad+\yVS\,\psi_{1644}^\dagger\beta_{1644}P_{1644}^{L}\psi_{1644}=0 .
\label{eq:fieldeqs}
\end{align}
%%</1644>
For $\PhiVS=\vVS/\sqrt2$ the first equation is the Dirac equation of Sec.~\ref{sec:dirac} with the
mass (\ref{eq:masses}). The second shows that the fermions act as a source of the Vorton-Scalar: a
medium of fermions, with $\langle\bar\psi P^{L}\psi\rangle\neq0$, displaces the scalar field from its
vacuum value. This is what makes the fermion mass depend on the density (Sec.~\ref{sec:matter}).

\subsection{The internal symmetry of the zero-mode sector}
\label{sec:internal}

The pairing calculation of the following sections uses fields that do not depend on the four extra
coordinates: all momenta lie in ordinary space. In this zero-mode sector the dynamics is fixed by five
Hermitian matrices only, $\alpha^1,\alpha^2,\alpha^3$, $\beta$ and $i\beta\chi$, and an \emph{internal
symmetry} is a unitary transformation $\psi\to e^{i\theta a}\psi$ of the sixteen components whose
generator $a$ commutes with all five. The set of all such generators was computed from the matrices of
each spacetime, by solving the linear equations that express the commutation.
Table~\ref{tab:symmetry} gives the result.

\begin{table*}
\caption{\label{tab:symmetry}The internal symmetry of the models in the sector of fields that do not
depend on the extra coordinates, computed from the Dirac matrices of the two spacetimes. The
\emph{commutant} is the set of all matrices that commute with the matrices that fix the dynamics;
its dimension is the number of independent symmetry generators. $P$ is the traceless central
generator. All entries have the values required by the structure
$su(2)_A\oplus su(2)_B\oplus u(1)_F\oplus u(1)_E$ described in the text. A number that the program
prints as a decimal fraction and that agrees with an integer in all printed digits is shown as that
integer.}
\begin{ruledtabular}
\begin{tabular}{>{\raggedright\arraybackslash}p{0.56\textwidth}>{\raggedright\arraybackslash}p{0.19\textwidth}>{\raggedright\arraybackslash}p{0.19\textwidth}}
Quantity computed from the matrices & $7+1$ & $4+4$ \\
\hline
{{ROWS:cjt_symmetry_table|rows=0-3,5,8,11,13,14,16-23}}
\end{tabular}
\end{ruledtabular}
\end{table*}

The result is \emph{the same for both spacetimes}. There are eight independent generators, and they
form the Lie algebra $su(2)_A\oplus su(2)_B\oplus u(1)_F\oplus u(1)_E$. Here $u(1)_F$ is the fermion
number, and $u(1)_E$ is generated by a matrix $P$ with $P^2=1$ that distinguishes the doublets of
$SU(2)_A$, which have $P=+1$, from those of $SU(2)_B$, which have $P=-1$. The eight states of a
particle of given momentum form the representation $(2,1)\oplus(1,2)$. The group
$[SU(2)_A\times SU(2)_B]/Z_2$ is $SO(4)$, and the four vector pair operators $\Cm\gamma^a$ with
$a\in\{0,5,6,7\}$ form its four-vector, the representation $(2,2)$, with fermion number 2. There is
exactly one invariant quadratic form on these four operators. It is
\begin{equation}
\mathrm{diag}(+,+,+,+)\ \ \text{in }7+1,\qquad\mathrm{diag}(+,-,-,-)\ \ \text{in }4+4,
\label{eq:invform}
\end{equation}
which is the metric of each spacetime on the four extra axes, up to an overall sign. It reappears in
the energy of the superfluid, Eq.~(\ref{eq:quartic}).

The two spacetimes differ in \emph{what the six generators of $SO(4)$ are}.
%%<1671>
On the $7+1$ spacetime they are the matrices $i\gamma^a\gamma^b$ with $a,b\in\{0,5,6,7\}$: the
rotations of the four extra space axes. They are a symmetry of the full theory, in which they rotate
the coordinates as well. If the fields may depend on $x0$, the rotations that involve the hidden axis
are no longer internal, and four generators remain, those of the rotations of $x5,x6,x7$ and the
fermion number; if the fields depend on all seven space coordinates only the fermion number remains
(Table~\ref{tab:symmetry}). On the $4+4$ spacetime of the companion paper the same algebra appears,
but three of its generators are not transformations of the spacetime at all \cite{Companion1644}.
%%</1671>
%%<1644>
Three of them are the rotations $i\gamma^s\gamma^t$ of the three additional times, $s,t\in\{5,6,7\}$.
The other three are the matrices $\gamma^0\gamma^t$. On the $4+4$ spacetime these matrices are
\emph{Hermitian} (last rows of Table~\ref{tab:symmetry}), so that $e^{i\theta\gamma^0\gamma^t}$ is
unitary. It is \emph{not} the Lorentz boost $e^{\theta\gamma^0\gamma^t/2}$ that mixes $x0$ with a time
coordinate, for a boost is not unitary: the algebra of the boosts and the algebra found here are two
different real forms, $\so(1,3)$ and $\so(4)$, of the same complex algebra. The $SO(4)$ of the $4+4$
model is therefore a \emph{hidden symmetry}. It acts on the components of the field, has no
counterpart among the transformations of the coordinates, and exists only as long as the fields do not
depend on $x0$: with such a dependence four generators remain, those of the rotations of the
additional times and the fermion number (second row of Table~\ref{tab:symmetry}). On the $7+1$
spacetime of the companion paper the same $SO(4)$ is the rotation group of four extra space axes, a
symmetry of the spacetime itself \cite{Companion1671}. This difference decides the fate of the Higgs
hypothesis (Sec.~\ref{sec:higgs}).
%%</1644>

When the chemical potential vanishes, the zero-mode sector has still more symmetry: further
generators exist that mix the field with its conjugate, and the full group is $O(4)\times O(4)$, of
dimension 12 (Table~\ref{tab:symmetry}). The reason is that the sixteen components are equivalent to
four ordinary Dirac fields of equal mass. A chemical potential for the fermion number removes these
generators. This matters for the vacuum (Sec.~\ref{sec:vacuum}).

\subsection{Vortices and vortons}
\label{sec:vortons}

This subsection collects known results that apply to the Vorton-Scalar. None of them is computed in
this paper, and each is attributed. Two kinds of vortex must be kept apart.

\emph{Vortices of the superfluid.} The order parameter of the superfluid has a phase, the Goldstone
field of the fermion number. A line around which this phase winds by $2\pi$ is a quantized vortex of
the superfluid, with circulation $2\pi/(2m)$ per unit of winding, since the pair has the mass $2m$
\cite{Onsager1949,Feynman1955}. In its core the condensate vanishes, and the scalar density of the
fermions differs from its value in the bulk; by the second of Eqs.~(\ref{eq:fieldeqs}) so does the
Vorton-Scalar. Every such vortex therefore carries a tube of displaced Vorton-Scalar field along with
it. This is the sense in which the scalar was named.

\emph{Vortices of the Vorton-Scalar itself.} A complex scalar with the potential of
Eq.~(\ref{eq:LVS}) has static solutions in which the phase winds by $2\pi n$ around a line and
$|\PhiVS|$ vanishes on the line \cite{NielsenOlesen1973}. For a global symmetry such as $\UVS$ the
energy per length grows logarithmically with the distance from the line. A fermion that obtains its
mass from the scalar through a Yukawa coupling is massless on the line, and the Dirac equation in the
background of the vortex has $|n|$ normalizable solutions of zero energy that move along the line at
the speed of light \cite{JackiwRossi1981}; the first of Eqs.~(\ref{eq:fieldeqs}) is of exactly this
type. Such zero modes make the line a perfect conductor of the corresponding charge, here the fermion
number \cite{Witten1985}. A closed loop that carries such a current has a conserved number of charge
carriers and a conserved winding of the current, and the resulting angular momentum can balance the
tension of the loop and prevent its collapse. Such stationary loops are \emph{vortons}
\cite{DavisShellard1989}. Whether they are stable, and with what abundance they would be produced,
depends on details that have been studied for other models and are open here.

\emph{The hidden ring.} In the second stage of the program the coordinate $x0$ curls up into a ring of
radius $R$. A Vorton-Scalar whose phase winds around that ring, $\PhiVS\propto e^{i\,n\,x0/R}$, is a
configuration with a constant neutral current along the hidden axis, of magnitude $n\vVS^2/R$. By
Eq.~(\ref{eq:LNC}) it acts on the fermions as a constant axial potential along $x0$. This is a natural
object of study for the second stage; nothing about it is claimed here.

\subsection{What has been established, and what has not}

%%<1671>
Established: a field with the required properties exists and is defined by the Lagrangian
(\ref{eq:LVS}). It is a complex scalar whose neutral current couples to the axial current of the
fermions with a constant of the Fermi type, Eq.~(\ref{eq:LNC}), and from which the mass of the
fermions originates, Eq.~(\ref{eq:masses}). The internal symmetry of the zero-mode sector is
$SU(2)_A\times SU(2)_B\times U(1)_F\times U(1)_E$, the rotations of the four extra space axes together
with two phases (Table~\ref{tab:symmetry}).

Not established: the completion of the theory at high energy, since the couplings have negative mass
dimension in eight dimensions; the stability and abundance of vortons; and whether nature contains
this field.
%%</1671>
%%<1644>
Established: in the fully covariant $4+4$ theory the construction is impossible
(Sec.~\ref{sec:nogo}); in the quantized sector it exists and is defined by the Lagrangian
(\ref{eq:LVS}). There the Vorton-Scalar is a complex scalar whose neutral current couples to the
axial current of the fermions with a constant of the Fermi type, Eq.~(\ref{eq:LNC}), and from which
the mass of the fermions originates, Eq.~(\ref{eq:masses}). The internal symmetry of the zero-mode
sector is $SU(2)_A\times SU(2)_B\times U(1)_F\times U(1)_E$, of which the $SO(4)$ is a hidden symmetry
(Table~\ref{tab:symmetry}).

Not established: the consistency of the model with the full $4+4$ spacetime, since the restriction to
fields independent of the additional times is an assumption; the completion of the theory at high
energy, since the couplings have negative mass dimension in eight dimensions; the stability and
abundance of vortons; and whether nature contains this field.
%%</1644>
